

\documentclass{article}
\usepackage{pathfinder-note}

\title{Mechanism Design for Multi-Agent Credit Assignment: Q$\times$P}

\begin{document}

\maketitle

\section*{The Pair}

Paper~Q develops a mechanism design framework for AI agents whose alignment (preferences) and capabilities (feasible actions, information) are unknown. The central insight is a \emph{one-sided imitation structure}: capabilities can be concealed but not counterfeited. This asymmetry yields a revelation principle, a characterization of implementable policies via nested cyclical monotonicity, and conditions under which eliciting higher-order beliefs can discipline multiple agents. Applications include sandbagging, peer scoring, and reward shaping.

Paper~P proposes MARS-RA, which reformulates credit assignment in cooperative multi-agent reinforcement learning as a rank aggregation problem. Contribution-based pairwise comparisons are generated by large multimodal models and converted into potential-based reward shaping. The framework assumes cooperative agents and honest LMM-generated comparisons.

\section*{Connexion Pursued}

The hypothesis connecting Q and P asks whether Q's mechanism design tools can harden P's credit assignment against strategic manipulation---particularly sandbagging, where a more capable agent pretends to be less capable to receive easier tasks or avoid scrutiny. The ledger identifies a critical tension: P assumes cooperative MARL with honest comparisons, while Q explicitly addresses strategic agents who may conceal capabilities \ledger{1}.

The connexion rests on whether agents in P's setting have incentives to manipulate the LMM-generated comparisons that drive reward shaping. If agents can strategically influence how their contributions are perceived or compared, Q's mechanisms for eliciting honest revelation of capabilities become relevant.

\section*{Supported Findings}

The ledger establishes that Q's applications explicitly include sandbagging, peer scoring, and reward shaping---all structurally present in P's framework \ledger{2}. This overlap is not coincidental: both papers address multi-agent settings where individual contributions must be assessed and rewarded.

The one-sided imitation structure from Q provides a theoretical lens: if agents in P's embodied MARL setting can conceal capabilities (e.g., by underperforming to appear less competent) but cannot counterfeit capabilities they lack, then Q's revelation principle may apply. This would enable mechanisms that incentivize honest capability revelation alongside honest effort.

Q's treatment of higher-order beliefs for disciplining multiple agents suggests a path toward peer-scoring mechanisms that could complement or replace P's LMM-generated comparisons when strategic behavior is suspected.

\section*{Objections}

The primary objection is that P's cooperative MARL setting may not exhibit the strategic incentives Q's mechanisms are designed to address. If agents share aligned objectives and the LMM evaluator is trusted, the added complexity of Q's mechanism design may be unnecessary.

A second objection concerns the evaluator: P relies on LMM-generated comparisons, which introduces a different vulnerability---manipulation of the evaluator's perceptions rather than direct capability concealment. Q's framework assumes the mechanism designer observes actions but not types; if the LMM itself can be strategically influenced, the problem shifts from mechanism design to evaluator robustness.

\section*{Unresolved Issues}

Evidence is thin on whether embodied MARL agents in P's experimental tasks actually exhibit strategic manipulation. The ledger notes the hypothesis ``needs sharpening'' regarding when agents would have incentives to manipulate comparisons \ledger{2}. Without empirical evidence of sandbagging or strategic behavior in P's setting, the applicability of Q's mechanisms remains speculative.

The ledger does not establish whether Q's higher-order belief elicitation can be operationalized in P's dynamic, partially-observable embodied environments. The theoretical conditions for disciplining multiple agents assume information structures that may not hold in practice.

\section*{Smallest Next Step}

Examine P's experimental results and task specifications to identify whether any agent behaviors resemble strategic capability concealment. If none exist, determine whether the task incentives would permit sandbagging even if agents were strategically motivated. This grounds the theoretical connexion in empirical reality before designing Q-inspired mechanisms.

\end{document}